\subsubsection{矩阵树定理}

\textbf{General Form:}
The sum of weights of all spanning trees is given by the cofactor of the Laplacian matrix:
\[ \sum_{T \in \mathcal{T}} \prod_{e \in T} w_e = \det(L^*) \]

\textbf{1. Undirected Graph}
\[
L_{ij} = \begin{cases} 
\deg(i) & i = j \\ 
-w_{ij} & i \neq j 
\end{cases}
\]
\textbf{Result:} $\det(L)$ after removing \textbf{any} $k$-th row and column.

\textbf{2. Directed Graph (Root $r$)}
Remove the $r$-th row and column from $L$ to calculate determinant.

\noindent \textbf{A. Out-Tree ($r \to \text{all}$):} Use \textbf{In-Degree} \& $A^T$.
\[
L_{ij} = \begin{cases} 
\deg^{in}(i) & i = j \\ 
-w_{ji} & i \neq j \quad (\text{edge } j \to i)
\end{cases}
\]

\noindent \textbf{B. In-Tree ($\text{all} \to r$):} Use \textbf{Out-Degree} \& $A$.
\[
L_{ij} = \begin{cases} 
\deg^{out}(i) & i = j \\ 
-w_{ij} & i \neq j \quad (\text{edge } i \to j)
\end{cases}
\]
当然，此处的 $w$ 可以替换为任意有交换律的东西，比如形式幂级数。

\begin{lstlisting}
void solve(){
    n=d,m=d;ll o=d;
    f(i,1,n)f(j,1,n)a[i][j]=0;//0 无向；1 有向且 1 为根外向树
    if(o==0){
        f(i,1,m){
            ll u=d,v=d,w=d;
            a[u][u]+=w,a[v][v]+=w;
            a[u][v]-=w,a[v][u]-=w;
        }
        f(i,2,n)f(j,2,n)a[i-1][j-1]=a[i][j];
        printf("%lld\n",det(n-1,a));
    }if(o==1){
        f(i,1,m){
            ll u=d,v=d,w=d;
            a[v][v]+=w;
            a[v][u]-=w;
        }
        f(i,2,n)f(j,2,n)a[i-1][j-1]=a[i][j];
        printf("%lld\n",det(n-1,a));
    }
}
\end{lstlisting}

\paragraph{使用条件与多测} 每组必须先清整个当前 $n\times n$ 矩阵，已在 solve 开头补上；自环对生成树无贡献，建矩阵时应忽略（代码的加减会抵消）。平行边权相加；加权度数是边权和而非边数。原矩阵压缩与 det 均破坏输入；多次删不同根时须从副本构造。权值要处于交换环中，仅“有交换律”不够；若调用前文 Fermat 消元，仍要求素数模，形式幂级数不能直接套除法实现。$n=1$ 答案为 1。
